% Converted from E-noir.md.
% Required packages: \usepackage{amsmath,longtable,booktabs,array}

\section{The Noir Circuit}\label{the-noir-circuit}

We wrote the Noir circuit directly from the clearing constraints. It supplies an independent implementation of the rule and the general-purpose benchmark in Section 5.5.

\textbf{Statement.} The witness consists of four arrays of \(N\) unsigned 32-bit volumes and a salt. The crossing circuit publishes nine field elements (288 bytes) for the book digest and eight receipt values. The no-cross circuit publishes only the digest.

\textbf{Circuit model.} Noir constrains the result of ordinary arithmetic, including the sums used to compute depth. Its unsigned integer types enforce bounds, so a subtraction that would underflow cannot satisfy the circuit. An unconstrained function can supply a witness value, but the circuit must still check that value. The backend implements the polynomial domain, quotient, openings, and lookups.

\textbf{Executable volume.} An unconstrained function proposes the minimum of the two depths. The circuit subtracts this value from each depth in unsigned 64-bit arithmetic, enforcing the differences and their non-negativity (F9, F10, F12, F13). It then checks that their product is zero (F11). Together the checks force the proposed value to equal the smaller depth.

\textbf{Types.} Table 16 maps Noir types to the constraints they enforce. Arrays have exactly \(N\) entries. Native comparisons implement the inequalities in F14 and F19 without an explicit slack column or limb witnesses at the circuit-language level.

\textbf{Table 16.} Types of the circuit and the constraints they replace.

\begin{longtable}[]{@{}lll@{}}
\toprule\noalign{}
Quantity & Type & Replaces \\
\midrule\noalign{}
\endhead
\bottomrule\noalign{}
\endlastfoot
Volume at one tick & 32-bit unsigned & F1, F2, F24, F25 \\
Depths, \(Min\), surpluses, \(\Delta\) & 64-bit unsigned & F12, F13, and the range of every accumulated value \\
\(c\), \(d\), \(p^*\) & 32-bit unsigned & \\
Side & Boolean & \\
Array of length \(N\) & & F3, F4, F26, F27 \\
\end{longtable}

\textbf{Receipt.} An unconstrained function computes a candidate receipt from the book. The circuit then applies the scalar and per-tick checks in Table 17. By Theorem 1, every accepted candidate equals the rule's receipt. The backend handles reads at \(c\), \(d\), \(p^*\), and their neighbours as memory lookups because those positions depend on the witness.

\textbf{Table 17.} Steps of the clearing check in the circuit.

\begin{longtable}[]{@{}lll@{}}
\toprule\noalign{}
Step & Assertion & Constraints \\
\midrule\noalign{}
\endhead
\bottomrule\noalign{}
\endlastfoot
Receipt & \(c \le p^* \le d < N\); \(V_{max} \ge 1\); \(V_{min\Delta} \le P\); \(s\), \(P\), \(J\) zero when balanced & Scalar checks \\
Depth & Running sums from the top tick and from the bottom tick & F5--F8, F28--F31 \\
Each tick & The minimum; \(\Delta = SurpA + SurpB\); \(Min \le V_{max}\); \(\Delta \ge V_{min\Delta}\) on \([c,d]\) & F9--F14, F18, F19 \\
Plateau & \(Min[c] = Min[d] = V_{max}\); \(Min\) strictly lower at \(c - 1\) and \(d + 1\) where they exist & F15--F17 \\
Clearing tick & The surpluses at \(p^*\) are \(V_{min\Delta}\) on the long side and 0 on the other & F20 \\
Tie-break & If \(p^* > c\) and supply is long or the batch is balanced, \(\Delta[p^* - 1] > V_{min\Delta}\) & F21 \\
Rationing & \(P\) is the long side\textquotesingle s volume at \(p^*\); \(J\) is the short side\textquotesingle s eligible IO depth at \(p^*\) & F32, F33 \\
\end{longtable}

\textbf{Digest.} We pack the four volumes at each tick into \(b+2^{32}a+2^{64}b^{io}+2^{96}a^{io}\). A width-four Poseidon2 sponge with rate three and capacity one absorbs the resulting field elements. We initialize the state to \((\text{salt},0,0,N)\), permute once before absorption, and take the first state element as the output after absorption. The private salt prevents an observer from testing a guessed low-volume book against the digest without also guessing the salt. At 1,000 ticks a Pedersen hash of the unpacked volumes required 1.78 million gates. Our sponge uses about 25,000. This arrangement of the standard permutation has not received a separate security analysis.

\textbf{Size.} The circuit has 12,757 gates at 100 ticks, 100,357 at 1,000, and 197,666 at 2,000, about 100 per tick. The proof is 14,656 bytes and the verification key 3,680 bytes at every size. The number of ticks is a compile-time constant, so each book size has its own verification key.

\textbf{Tests.} Four of 20 library tests check hand-computed receipts for a demand-long book, a supply-long book, a balanced book, and a book with one maximizing tick and IO interest. Twelve tests submit false receipts and check which constraint rejects them. The remaining tests cover no-cross and the dependence of the digest on the book and salt. The Noir and Rust receipts agree on all 57 calibrated and replayed benchmark books and on the 362 books constructed from Cboe outcomes in Section 5.7.

\textbf{Toolchain.} We compile with \texttt{nargo} 1.0.0 and prove with Barretenberg 6.0.0 in its zero-knowledge UltraHonk mode with a Poseidon2 transcript. The circuit is 208 lines and its tests 161.
